%==============================================================================
% Chapter 3 -- Methods
%==============================================================================
\chapter{Methods}
\label{ch:methods}

% -----------------------------------------------------------------------------
% Adapted from Section 3 of the Design-CP paper. This is the technical core of
% the report and (as in the example) the longest chapter. Verify every numerical
% constant and symbol against the paper before submission; where the paper's
% appendices give more detail (device meshes, communication volumes), summarise
% here and cite the appendix.
% -----------------------------------------------------------------------------

As discussed in Chapter~\ref{ch:current-memory-reduction-methods}, the current barrier for designing
whole protein assembly de-novo is a memory problem. An all-atom
generative model can express a whole assembly, but cannot hold its quadratic
activations needed for pair-biased attention on one accelerator. This chapter introduces \emph{Design-CP}, two
context-parallel inference strategies that shard those activations across a
multi-GPU mesh while leaving the pretrained model weights untouched. Neither
scheme alters the computation the model performs: both reproduce single-GPU
inference numerically, and only the placement of the activations across devices
differs.

\section{RFdiffusion3 preliminaries and notation}
\label{sec:rfd3-prelim}

We build on RFdiffusion3 (RFD3) \parencite{butcher_novo_2025}, an all-atom denoising-diffusion model whose
diffusion module follows the AlphaFold3 design
\parencite{abramsonAccurateStructurePrediction2024} and which is
built on the AtomWorks framework \parencite{corley_accelerating_2025}. At each denoising step RFD3 jointly represents a
structure at two resolutions. A \emph{token} track carries a single
representation $S \in \mathbb{R}^{I \times c_s}$ and a pair representation
$Z \in \mathbb{R}^{I \times I \times c_z}$ over $I$ tokens; an \emph{atom} track
carries a single representation $A \in \mathbb{R}^{L \times c_\text{atom}}$ and a pair representation
$P \in \mathbb{R}^{L \times L \times c_\text{atompair}}$ over $L$ atoms. Attention
on each track adds a learned pair bias to the logits, e.g.
%
\begin{equation}
  \mathrm{Attn}(S) = \mathrm{softmax}\!\left(\frac{QK^\top}{\sqrt{d}} + B(Z)\right)V,
  \label{eq:pair-bias-attn}
\end{equation}
%
where $B(Z) \in \mathbb{R}^{I \times I \times H}$ projects the pair track into a
bias for each of the $H$ attention heads. Atom attention can be made \emph{sparse}
using a $k$-nearest-neighbour neighbourhood with $k \ll L$, so its compute cost is
$\mathcal{O}(Lk)$ rather than $\mathcal{O}(L^2)$. This is an approximation with respect to 
how the model was trained (which used the full matrix $P$), and when not used the pair storage $P$
remains quadratic in $L$. The single-GPU memory footprint is
therefore dominated by the $\mathcal{O}(I^2)$ and possibly the $\mathcal{O}(L^2)$ pair tensors,
which is what limits the size of assembly that can be modelled at once.
Sampling follows the EDM framework \parencite{karras_elucidating_2022} with the
AlphaFold3 schedule. We leave the stock RFD3 settings untouched:
$T=200$ denoising steps, two recycling iterations per step, and a sparse atom
neighbourhood of $k=128$ keys.

When designing with a pre-defined point symmetry, the diffusion model is always
run on the entire complex at every denoising step. For a configurable fraction of
the trajectory (default $90\%$), RFD3 then \emph{resymmetrises} its prediction by
extracting the coordinates of a single asymmetric subunit (ASU) and generating the
remaining copies by applying the point-group operations. The resulting
resymmetrised complex is the state that is fed into the next iteration of the
sampling loop. Because the network still denoises every atom of the assembly, it
can account for the full set of inter-subunit interactions, while the symmetry
constraint couples atom positions across subunits and thereby reduces the
effective degrees of freedom of the design to those of a single ASU. The per-step
procedure is set out in Appendix~\ref{app:implementation}.

\begin{figure}[t]
  \centering
  \includegraphics[width=\textwidth]{Figure_1.png}
  \caption{\textbf{Symmetric design with Design-CP.}
    \textbf{(a)} Schematic depiction of the sharding techniques implemented in
    Design-CP. Every square represents a sub-tensor of the self-attention matrix,
    and its colour represents its assigned device. 1D sharding partitions the
    queries across different GPUs, where they are used to calculate
    cross-attention against all keys. 2D sharding partitions both queries and keys
    and uses ring attention to compute attention scores on the fly.
    \textbf{(b)} Qualitative comparison of large-design samples. Designs of a
    $10{,}800$ amino-acid protein using different symmetries. Designs that exceed
    the native crop limits can show visible degradation when sampled without
    additional structure constraints; imposing a strong symmetry prior mitigates
    this effect by reducing the effective design space.}
  \label{fig:large-designs}
\end{figure}

\section{Design-CP 1D row-sharding}
\label{sec:cp-1d}

Our first scheme, implemented in PyTorch with NCCL and without custom kernels or
\texttt{DTensor} machinery, partitions the query dimension of every $I\times I$
and $L\times L$ operation across $P$ GPUs (Figure~\ref{fig:large-designs}a, center).
GPU~$p$ materialises only its row stripe of the pair track:
%
\begin{equation}
  \mathbf{Z}^{(p)} \in \mathbb{R}^{I_p \times I \times c_z}, \qquad
  I_p = \lfloor I/P \rfloor + \mathbb{I}[\,p < I \bmod P\,],
  \label{eq:row-stripe}
\end{equation}
%
where $\mathbb{I}[\,\cdot\,]\in\{0,1\}$ denotes the indicator function of its
predicate. When $I$ is not divisible by $P$, the floor $\lfloor I/P\rfloor$ leaves
a remainder of $I \bmod P$ elements that must still be assigned. We absorb this
remainder by giving the first $I \bmod P$ ranks one additional row each: rank $p$
receives the extra row precisely when $p < I \bmod P$, which is what the indicator
encodes. By construction $\sum_{p}I_p = I$, no element is dropped, and chunk sizes
differ by at most one across GPUs, so the load imbalance per attention block is
bounded by a single row regardless of $P$. The pair tracks $\mathbf{Z}^{(p)}$ and
the self-conditioning distogram $\mathbf{D}^{(p)}_\text{self}$ are in this way
\emph{never} gathered to their full $[I,I]$ shape.

Every self-attention block is reformulated as a cross-attention: GPU~$p$ projects
queries from its stripe $\mathbf{Q}^{(p)}=f_Q(\mathbf{S}^{(p)})\in
\mathbb{R}^{I_p\times H\times d}$, while keys and values are projected from the
\emph{full} single-track $\mathbf{S}$, which is replicated on every GPU. The pair
bias is drawn from the local stripe, $\mathbf{B}^{(p)}=f_B(\mathbf{Z}^{(p)})\in
\mathbb{R}^{I_p\times I\times H}$, and
%
\begin{equation}
  \operatorname{Attn}^{(p)} = \operatorname{softmax}\!\left(
    \mathbf{Q}^{(p)} \mathbf{K}^\top/\sqrt{d} + \mathbf{B}^{(p)}
  \right) \mathbf{V}.
  \label{eq:cross-attn}
\end{equation}
%
A single \textsc{AllGather} over the 1D track reconstructs
$\mathbf{S}=\bigoplus_{p}\mathbf{S}^{(p)}\in\mathbb{R}^{I\times c_s}$ after each
block, so that all GPUs hold identical $\mathbf{K}$ and $\mathbf{V}$ for the next
block. The same partition applies at the atom level. The dense atom pair track is
striped as $\mathbf{P}^{(p)}\in\mathbb{R}^{L_p\times L\times c_\text{atompair}}$,
reducing its per-GPU storage from $\mathcal{O}(L^{2})$ to
$\mathcal{O}(L^{2}/P)$. The sparse $k$-nearest-neighbour attention is then applied
per shard: each GPU gathers, from its stripe of $\mathbf{P}^{(p)}$, the $k$
neighbour entries for each of its $L_p$ query atoms, yielding a local bias
$\mathbf{P}^{(p)}_\text{sparse}\in\mathbb{R}^{L_p\times k\times c_\text{atompair}}$
at attention-computation cost $\mathcal{O}(L\,k/P)$. At diffusion-step boundaries,
rank~$0$ broadcasts the noised coordinates, sampled Gaussian noise and denoised
prediction so that all ranks share a single stochastic trajectory. Per-GPU
pair-track memory is therefore $\mathcal{O}(I^{2}/P)$ for $\mathbf{Z}$ and
$\mathcal{O}(L^{2}/P)$ for $\mathbf{P}$. A more detailed description is provided in
Appendix~\ref{app:implementation}.

\section{Design-CP 2D grid context parallelism}
\label{sec:cp-2d}

Our second scheme adopts the Fold-CP framework \parencite{lin_fold-cp_2026} and
specialises it to RFD3 (Figure~\ref{fig:large-designs}a, right). We arrange $P$ GPUs on a
$\sqrt{P}\times\sqrt{P}$ grid and tile the pair tensor into local quadrants
$\mathbf{Z}^{(r,c)}\in\mathbb{R}^{I_r\times I_c\times c_z}$ at grid position
$(r,c)$, with $I_r=I/\sqrt{P}$. Queries are sharded along the row axis and
replicated along the column axis; keys and values are sharded along the column
axis and circulated by $\sqrt{P}$ ring shifts, with an initial
$(r,c)\leftrightarrow(c,r)$ transpose to align them with the query rows. At each
ring step, a GPU computes attention between its resident $\mathbf{Q}$ rows and the
currently visited $\mathbf{K}$/$\mathbf{V}$ shard, using the local bias quadrant,
and merges the partial output into a running total via an online softmax
\parencite{milakov_online_2018,liu_ring_2023}. All tensors are represented as
PyTorch \texttt{DTensor}s; parameters are replicated across the grid.

Asymptotically, per-device pair-track memory is $\mathcal{O}(I^{2}/P)$, the same
as the 1D scheme, while $\mathbf{K}$ and $\mathbf{V}$ are ring-rotated rather than
replicated, so each device only ever holds an $\mathcal{O}(I/\sqrt{P})$ slab of
$\mathbf{K}$/$\mathbf{V}$ at a time. The $\sqrt{P}$ ring shifts per attention block
are overlapped with local computation.

Relative to Fold-CP, the RFD3 adaptation mainly concerns the sampling loop and the
atom-level path: the ring primitive is invoked inside every recycling iteration of
every denoising step, with rank-$0$ broadcasts at step boundaries that mirror the
1D scheme, and the sparse atom attention requires a distributed $k$-nearest-neighbour
search that avoids materialising any $[L,L]$ distance tensor. Concrete descriptions
of these adaptations are deferred to Appendix~\ref{app:implementation}. This 2D
grid scheme applies only when the available GPU count $P$ is a perfect square, so
that devices can be arranged as a $\sqrt{P}\times\sqrt{P}$ mesh.

%NOTE: Keep this commented out for now. I don't know if I want it 
% \section{Composition of symmetry with context parallelism}
% \label{sec:symmetry-cp}

% In both schemes the symmetrisation of Section~\ref{sec:rfd3-symmetry} is executed
% redundantly on every rank rather than being distributed: each device applies the
% point-group operations to full-size coordinates and obtains the same result, so
% every rank follows an identical sampling trajectory. The two schemes differ in how
% those identical inputs are produced.

% The 1D scheme keeps the coordinate state replicated across ranks. Its sampling
% loop broadcasts the sampled Gaussian noise and the network's clean prediction from
% rank~$0$ at every denoising step, and the initial coordinates once at the start;
% these broadcasts are a determinism guard that is present whether or not symmetry
% is active. The symmetrised path therefore consumes bit-identical inputs on every
% device, and the deterministic ASU extraction and frame application reproduce the
% same symmetrised coordinates everywhere, without any communication attributable to
% symmetry itself.

% The 2D scheme instead holds the coordinates as \texttt{DTensor}s sharded along the
% atom dimension, and obtains identical noise by drawing the full-shape tensor from
% per-rank RNG states held in lockstep before sharding it, rather than by
% broadcasting. Because the ASU is selected by boolean indexing over the atom
% dimension, the coordinates must be gathered before the point-group operations are
% applied and re-sharded afterwards, once for the noised coordinates and once for
% the prediction. Symmetric design does therefore add collectives in this scheme,
% but they act on the $\mathcal{O}(L)$ coordinate array rather than on the
% $\mathcal{O}(I^{2})$ pair track, and are negligible beside the $\sqrt{P}$ ring
% exchanges performed inside every attention block of every recycling iteration.

% Neither scheme alters the sharding of the pair track, the ring schedule or the
% model weights; symmetry acts only on the coordinate state between denoising steps.
% This is the operational sense in which the symmetrisation procedure is
% \emph{orthogonal} to context parallelism, and it is the reason a single set of
% pretrained weights designs both the unsymmetrised baselines and the icosahedral
% and octahedral assemblies of Chapter~\ref{ch:results}. Further detail is given in
% Appendix~\ref{app:implementation}.

\section{In-silico evaluation}
\label{sec:evaluation}

A standard \emph{de novo} design pipeline often couples a structure generator
(RFdiffusion or RFD3), a sequence designer such as ProteinMPNN
\parencite{dauparasRobustDeepLearning2022}, and an external structure predictor
for refolding-based validation. For very large multimeric assemblies, however,
prediction-based oracles become both more expensive and less reliable:
AlphaFold-Multimer accuracy, for example, has been found to degrade with chain
count \parencite{bryant_predicting_2022}. We therefore assess design quality
primarily through \emph{in silico} structural sanity checks computed directly from
the generated all-atom coordinates, and we organise the assessment into two
questions: (i)~does sampling at this scale via Design-CP preserve the per-chain
backbone quality of the original RFD3 model, and (ii)~does the symmetrised
assembly realise the intended interfaces, consistently with a functional natural
baseline.

The metrics fall into two families. \emph{Backbone sanity} metrics flag designs
whose monomeric chain is geometrically broken, irrespective of any symmetry
consideration: chain breaks, backbone heavy-atom clashes, the non-loop fraction
and the maximum C$\alpha$--C$\alpha$ bond deviation. \emph{Symmetric-interface
sanity} metrics use the full symmetrised complex and probe its inter-subunit
geometry: ASU clashes, mean contacts per interface, minimum inter-chain distance
and the number of interfaces that make contact. For icosahedral designs these are
computed between the ASU and its eight nearest icosahedral neighbours. Where a
natural reference particle is available we overlay it in every panel; for
icosahedral designs we use lumazine synthase from \emph{Aquifex aeolicus}
(PDB~1NQX, biological assembly~1: a 60-subunit $T{=}1$ icosahedral capsid)
\parencite{zhang_structure-based_2003}, selected as a naturally occurring
nanoparticle that has previously been employed as a vaccine scaffold
\parencite{ladenstein_second_2020,joseph_lumazine_2025}. Precise definitions, with
every threshold, are collected in Appendix~\ref{app:metric-defs}.
